\subsection{正交性。双线性或半双线性映射的直和。}

本节中，A 与 B 表示环，E 表示左 A-模，F 表示右（相应地左）B-模，G 表示 (A, B)-双模，Φ 表示 E × F 到 G 的双线性（相应地，关于 B 的一个给定反自同构 J 的半双线性）映射。

定义 6. --- 称两个元素 \(x \in E\) 与 \(y \in F\) 关于 \(\Phi\) 正交，如果 \(\Phi(x, y) = 0\) 。称两个子集 \(E' \subset E\) 与 \(F' \subset F\) 正交，如果对任意 \(x \in E'\) 与 \(y \in F'\) ，\(x\) 与 \(y\) 正交。E（相应地 F）中与 F 的给定子模 N（相应地 E 的子模 M）正交的元素所成之集是 E（相应地 F）的一个子模，称为与 N（相应地 M）完全正交（或简称正交）的子模，记作 \(N^0\) （相应地 \(M^0\) ）。

设 H 与 H′ 是 E 或 F 的两个子模。有 \(\mathrm{H}\subset(\mathrm{H}^{0})^{0}\) （记作 \(H^{00}\) ）；若 \(H\subset H'\) ，则 \(H^{\prime0}\supset H^{0}\) 。由此得 \(H^{0}\supset(H^{00})^{0}\) 且 \(H^{0}\subset(H^{0})^{00}\) ；令

\[
\mathrm{H} ^ {0 0 0} = (\mathrm{H} ^ {0 0}) ^ {0} = (\mathrm{H} ^ {0}) ^ {0 0} = ((\mathrm{H} ^ {0}) ^ {0}) ^ {0},
\]

便有 \(\mathbf{H}^0 = \mathbf{H}^{000}\) 。

映射 \(\Phi\) 退化（第 1 小节，定义 3）必须且只需两个子模 E \(^{0}\) 、F \(^{0}\) 中至少有一个 \(\neq \{0\}\) 。显然，\(\Phi(x, y)\) 在给 \(x\) （相应地 \(y\) ）加上

F \(^{0}\) （相应地 E \(^{0}\) ）中的一个元素时不改变，因而 Φ 通过过渡到商，在 (E/F \(^{0}\) ) × (F/E \(^{0}\) ) 上定义了一个双线性（或半双线性）映射；它显然非退化；称之为 Φ 的相伴非退化双线性（或半双线性）映射。

设 \((\mathrm{E}_{i})_{i\in\mathrm{I}}\) 为左 A-模的一个族，\((\mathrm{F}_{i})_{i\in\mathrm{I}}\) 为右 B-模（相应地左 B-模）的一个族，\(\Phi_{i}\) 为 \(E_{i}\times F_{i}\) 到 G 中的一个双线性（相应地关于 J 的右半双线性）映射。用 E（相应地 F）表示诸 \(E_{i}\) （相应地 \(F_{i}\) ）的直和模。我们立即看到，映射 \(\Phi\) （从 \(E\times F\) 到 G 中）由

\[
\Phi ((x _ {i}), (y _ {i})) = \sum_ {i} \Phi_ {i} (x _ {i}, y _ {i}) \quad (x _ {i} \in \mathrm{E} _ {i}, y _ {i} \in \mathrm{F} _ {i})\tag{10}
\]

（这个和有意义，因为除有限项外其各项均为零）定义，它是双线性的（相应地关于 J 是右半双线性的）。称它为诸映射 \(\Phi_{i}\) 的直和。显然，\(\mathbf{E}_{i}\) 与 \(\mathbf{F}_{j}\) 关于 \(\Phi\) 正交（\(i \neq j\) ）。反之，设 \(\Phi\) 为 \(\mathbf{E} \times \mathbf{F}\) 到 G 中的一个双线性或半双线性映射，并设 E 是子模 \((\mathbf{E}_{i})_{i \in I}\) 的直和，F 是子模 \((\mathbf{F}_{i})_{i \in I}\) 的直和，使得 \(\mathbf{E}_{i}\) 与 \(\mathbf{F}_{j}\) 正交（对 \(i \neq j\)）；则 \(\Phi\) 是它在积 \(\mathbf{E}_{i} \times \mathbf{F}_{i} (i \in I)\) 上的限制的直和。

\(\Phi\) 非退化必须且只需诸 \(\Phi_{i}\) 都非退化；在这些条件下，与 \(\mathbf{E}_{i}\) 正交的子模是 \(\sum_{j\neq i}\mathbf{F}_{j}\) 。

